% =====================================================================
%  Introduction : importée par main.tex avec % =====================================================================
%  Introduction : importée par main.tex avec % =====================================================================
%  Introduction : importée par main.tex avec % =====================================================================
%  Introduction : importée par main.tex avec \include{introduction}.
%  Section non numérotée, ajoutée au sommaire. Pas de figure flottante ici
%  (la numérotation des figures suit celle des chapitres) : le plan est un
%  dessin TikZ en place.
% =====================================================================

\phantomsection
\section*{Introduction}
\addcontentsline{toc}{section}{Introduction}

Ce document rassemble les notes d'un cours consacré à OpenStack, plateforme de
\gls{iaas} open source, et à ses principaux services. Il part d'une présentation générale
(cloud, virtualisation, architecture) puis consacre un chapitre à chaque service :
identité (Keystone), calcul (Nova, Placement), images (Glance), réseau (Neutron),
stockage bloc et objet (Cinder, Swift), tableau de bord (Horizon), répartition de charge
(Octavia), orchestration (Heat), bare metal (Ironic), Kubernetes managé (Magnum) et bases
de données (Trove). Un dernier chapitre présente les méthodes de déploiement.

\begin{center}
\begin{tikzpicture}[>=Stealth, font=\scriptsize,
    boite/.style={draw=insarouge, fill=insarouge!6, thick, rounded corners=2pt,
                  text width=3.35cm, minimum height=2.5cm, align=center},
    num/.style={draw=insarouge, fill=insarouge, text=white, font=\scriptsize\bfseries,
                rounded corners=2pt, minimum height=5.5mm, inner xsep=5pt}]
  \node[boite] (b1) at (0,0)
    {\textbf{Présentation générale}\\[3pt] cloud, IaaS, architecture d'OpenStack};
  \node[boite] (b2) at (4,0)
    {\textbf{Services de base}\\[3pt] Keystone, Nova, Placement, Glance, Neutron, Cinder, Horizon, Octavia};
  \node[boite] (b3) at (8,0)
    {\textbf{Services avancés}\\[3pt] Swift, Heat, Ironic, Magnum, Trove};
  \node[boite] (b4) at (12,0)
    {\textbf{Déploiement}\\[3pt] DevStack, Kolla-Ansible, installation manuelle};
  \node[num] at (0,1.55)  {chap.~\ref{chap:presentation}};
  \node[num] at (4,1.55)  {chap.~\ref{chap:keystone}--\ref{chap:octavia}};
  \node[num] at (8,1.55)  {chap.~\ref{chap:swift}--\ref{chap:trove}};
  \node[num] at (12,1.55) {chap.~\ref{chap:deploiement}};
  \foreach \i/\j in {1/2, 2/3, 3/4}
    \draw[->, insarouge, thick] (b\i.east) -- (b\j.west);
\end{tikzpicture}
\end{center}

Les termes techniques sont définis dans le glossaire, en fin de document ; la
documentation officielle d'OpenStack \parencite{openstack-docs} sert de référence commune
à tous les chapitres.

\paragraph{Prérequis et conventions.} Notions de Linux et de ligne de commande,
de réseau (adressage IP, VLAN, routage) et de virtualisation (machine virtuelle,
hyperviseur). Les exemples de commandes utilisent le client en ligne de commande
unifié \texttt{openstack}.
.
%  Section non numérotée, ajoutée au sommaire. Pas de figure flottante ici
%  (la numérotation des figures suit celle des chapitres) : le plan est un
%  dessin TikZ en place.
% =====================================================================

\phantomsection
\section*{Introduction}
\addcontentsline{toc}{section}{Introduction}

Ce document rassemble les notes d'un cours consacré à OpenStack, plateforme de
\gls{iaas} open source, et à ses principaux services. Il part d'une présentation générale
(cloud, virtualisation, architecture) puis consacre un chapitre à chaque service :
identité (Keystone), calcul (Nova, Placement), images (Glance), réseau (Neutron),
stockage bloc et objet (Cinder, Swift), tableau de bord (Horizon), répartition de charge
(Octavia), orchestration (Heat), bare metal (Ironic), Kubernetes managé (Magnum) et bases
de données (Trove). Un dernier chapitre présente les méthodes de déploiement.

\begin{center}
\begin{tikzpicture}[>=Stealth, font=\scriptsize,
    boite/.style={draw=insarouge, fill=insarouge!6, thick, rounded corners=2pt,
                  text width=3.35cm, minimum height=2.5cm, align=center},
    num/.style={draw=insarouge, fill=insarouge, text=white, font=\scriptsize\bfseries,
                rounded corners=2pt, minimum height=5.5mm, inner xsep=5pt}]
  \node[boite] (b1) at (0,0)
    {\textbf{Présentation générale}\\[3pt] cloud, IaaS, architecture d'OpenStack};
  \node[boite] (b2) at (4,0)
    {\textbf{Services de base}\\[3pt] Keystone, Nova, Placement, Glance, Neutron, Cinder, Horizon, Octavia};
  \node[boite] (b3) at (8,0)
    {\textbf{Services avancés}\\[3pt] Swift, Heat, Ironic, Magnum, Trove};
  \node[boite] (b4) at (12,0)
    {\textbf{Déploiement}\\[3pt] DevStack, Kolla-Ansible, installation manuelle};
  \node[num] at (0,1.55)  {chap.~\ref{chap:presentation}};
  \node[num] at (4,1.55)  {chap.~\ref{chap:keystone}--\ref{chap:octavia}};
  \node[num] at (8,1.55)  {chap.~\ref{chap:swift}--\ref{chap:trove}};
  \node[num] at (12,1.55) {chap.~\ref{chap:deploiement}};
  \foreach \i/\j in {1/2, 2/3, 3/4}
    \draw[->, insarouge, thick] (b\i.east) -- (b\j.west);
\end{tikzpicture}
\end{center}

Les termes techniques sont définis dans le glossaire, en fin de document ; la
documentation officielle d'OpenStack \parencite{openstack-docs} sert de référence commune
à tous les chapitres.

\paragraph{Prérequis et conventions.} Notions de Linux et de ligne de commande,
de réseau (adressage IP, VLAN, routage) et de virtualisation (machine virtuelle,
hyperviseur). Les exemples de commandes utilisent le client en ligne de commande
unifié \texttt{openstack}.
.
%  Section non numérotée, ajoutée au sommaire. Pas de figure flottante ici
%  (la numérotation des figures suit celle des chapitres) : le plan est un
%  dessin TikZ en place.
% =====================================================================

\phantomsection
\section*{Introduction}
\addcontentsline{toc}{section}{Introduction}

Ce document rassemble les notes d'un cours consacré à OpenStack, plateforme de
\gls{iaas} open source, et à ses principaux services. Il part d'une présentation générale
(cloud, virtualisation, architecture) puis consacre un chapitre à chaque service :
identité (Keystone), calcul (Nova, Placement), images (Glance), réseau (Neutron),
stockage bloc et objet (Cinder, Swift), tableau de bord (Horizon), répartition de charge
(Octavia), orchestration (Heat), bare metal (Ironic), Kubernetes managé (Magnum) et bases
de données (Trove). Un dernier chapitre présente les méthodes de déploiement.

\begin{center}
\begin{tikzpicture}[>=Stealth, font=\scriptsize,
    boite/.style={draw=insarouge, fill=insarouge!6, thick, rounded corners=2pt,
                  text width=3.35cm, minimum height=2.5cm, align=center},
    num/.style={draw=insarouge, fill=insarouge, text=white, font=\scriptsize\bfseries,
                rounded corners=2pt, minimum height=5.5mm, inner xsep=5pt}]
  \node[boite] (b1) at (0,0)
    {\textbf{Présentation générale}\\[3pt] cloud, IaaS, architecture d'OpenStack};
  \node[boite] (b2) at (4,0)
    {\textbf{Services de base}\\[3pt] Keystone, Nova, Placement, Glance, Neutron, Cinder, Horizon, Octavia};
  \node[boite] (b3) at (8,0)
    {\textbf{Services avancés}\\[3pt] Swift, Heat, Ironic, Magnum, Trove};
  \node[boite] (b4) at (12,0)
    {\textbf{Déploiement}\\[3pt] DevStack, Kolla-Ansible, installation manuelle};
  \node[num] at (0,1.55)  {chap.~\ref{chap:presentation}};
  \node[num] at (4,1.55)  {chap.~\ref{chap:keystone}--\ref{chap:octavia}};
  \node[num] at (8,1.55)  {chap.~\ref{chap:swift}--\ref{chap:trove}};
  \node[num] at (12,1.55) {chap.~\ref{chap:deploiement}};
  \foreach \i/\j in {1/2, 2/3, 3/4}
    \draw[->, insarouge, thick] (b\i.east) -- (b\j.west);
\end{tikzpicture}
\end{center}

Les termes techniques sont définis dans le glossaire, en fin de document ; la
documentation officielle d'OpenStack \parencite{openstack-docs} sert de référence commune
à tous les chapitres.

\paragraph{Prérequis et conventions.} Notions de Linux et de ligne de commande,
de réseau (adressage IP, VLAN, routage) et de virtualisation (machine virtuelle,
hyperviseur). Les exemples de commandes utilisent le client en ligne de commande
unifié \texttt{openstack}.
.
%  Section non numérotée, ajoutée au sommaire. Pas de figure flottante ici
%  (la numérotation des figures suit celle des chapitres) : le plan est un
%  dessin TikZ en place.
% =====================================================================

\phantomsection
\section*{Introduction}
\addcontentsline{toc}{section}{Introduction}

Ce document rassemble les notes d'un cours consacré à OpenStack, plateforme de
\gls{iaas} open source, et à ses principaux services. Il part d'une présentation générale
(cloud, virtualisation, architecture) puis consacre un chapitre à chaque service :
identité (Keystone), calcul (Nova, Placement), images (Glance), réseau (Neutron),
stockage bloc et objet (Cinder, Swift), tableau de bord (Horizon), répartition de charge
(Octavia), orchestration (Heat), bare metal (Ironic), Kubernetes managé (Magnum) et bases
de données (Trove). Un dernier chapitre présente les méthodes de déploiement.

\begin{center}
\begin{tikzpicture}[>=Stealth, font=\scriptsize,
    boite/.style={draw=insarouge, fill=insarouge!6, thick, rounded corners=2pt,
                  text width=3.35cm, minimum height=2.5cm, align=center},
    num/.style={draw=insarouge, fill=insarouge, text=white, font=\scriptsize\bfseries,
                rounded corners=2pt, minimum height=5.5mm, inner xsep=5pt}]
  \node[boite] (b1) at (0,0)
    {\textbf{Présentation générale}\\[3pt] cloud, IaaS, architecture d'OpenStack};
  \node[boite] (b2) at (4,0)
    {\textbf{Services de base}\\[3pt] Keystone, Nova, Placement, Glance, Neutron, Cinder, Horizon, Octavia};
  \node[boite] (b3) at (8,0)
    {\textbf{Services avancés}\\[3pt] Swift, Heat, Ironic, Magnum, Trove};
  \node[boite] (b4) at (12,0)
    {\textbf{Déploiement}\\[3pt] DevStack, Kolla-Ansible, installation manuelle};
  \node[num] at (0,1.55)  {chap.~\ref{chap:presentation}};
  \node[num] at (4,1.55)  {chap.~\ref{chap:keystone}--\ref{chap:octavia}};
  \node[num] at (8,1.55)  {chap.~\ref{chap:swift}--\ref{chap:trove}};
  \node[num] at (12,1.55) {chap.~\ref{chap:deploiement}};
  \foreach \i/\j in {1/2, 2/3, 3/4}
    \draw[->, insarouge, thick] (b\i.east) -- (b\j.west);
\end{tikzpicture}
\end{center}

Les termes techniques sont définis dans le glossaire, en fin de document ; la
documentation officielle d'OpenStack \parencite{openstack-docs} sert de référence commune
à tous les chapitres.

\paragraph{Prérequis et conventions.} Notions de Linux et de ligne de commande,
de réseau (adressage IP, VLAN, routage) et de virtualisation (machine virtuelle,
hyperviseur). Les exemples de commandes utilisent le client en ligne de commande
unifié \texttt{openstack}.
